\documentclass{newlayout}
%Bitte hier den enstprechenden Ort einsetzen z.B. Braunschweig und die Akademienummer
\Akademie{Ort}{2014}{1}

\usepackage[ngerman, english]{babel}
\usepackage{misc}
\usepackage{multicol}

%\usepackage{amsmath}%wird automatisch durch newlayout.cls geladen
\usepackage{amsfonts}
\usepackage{tikz}
\usepackage{blindtext}
\usepackage{lmodern}
\usepackage{url}
\def\UrlBreaks{\do\a\do\b\do\c\do\d\do\e\do\f\do\g\do\h\do\i\do\j\do\k\do\l%
\do\m\do\n\do\o\do\p\do\q\do\r\do\s\do\t\do\u\do\v\do\w\do\x\do\y\do\z\do\0%
\do\1\do\2\do\3\do\4\do\5\do\6\do\7\do\8\do\9\do\-\do\_\do\/\do\%}
\urlstyle{same}

% hinzugef�gt, um Fehler 'pdfTeX error (font expansion): auto expansion is only possible with scalable' zu vermeiden
\setkomafont{descriptionlabel}{\normalfont\bfseries}
\addtokomafont{paragraph}{\normalfont}
\usepackage{footnote}
\usepackage[flushmargin,hang,ragged]{footmisc}
\deffootnote{1em}{1em}{%
\textsuperscript{\thefootnotemark\ }
}
%\setlength{\abovedisplayskip}{5pt}
%\setlength{\belowdisplayskip}{5pt}


%%%%%Mathe-Definitionen
\newtheorem{Def}{Definition}
\newtheorem{Sat}{Satz}
\newtheorem{Bew}{Beweis}
\newtheorem{Thm}{Theorem}

\setlength\abovedisplayshortskip{0pt}
\setlength\belowdisplayshortskip{0pt}
\setlength\abovedisplayskip{3pt}
\setlength\belowdisplayskip{3pt}
%%%%Ende Mathe-Definitionen

\begin{document}

 %   \input{titel}
 \setcounter{page}{3}

\setcounter{tocdepth}{1}
 \tableofcontents

   \setcounter{secnumdepth}{1}


\setcounter{page}{7}
\setcounter{chapter}{0}

%Angabe, bis zu welcher Stufe die sections im Text nummeriert werden sollen.
      \settocdepth{2}



\course{1.2}{Lineare Algebra}%%% 
\begin{coursetitle}
  \centerline{Learning from Data aka i++} 
  \bigskip
  \Large \centerline{Mathematische Methoden des digitalen Zeitalters}
  \bigskip
\end{coursetitle}
%\begin{dsafigure}
%\begin{center}
%\includegraphics[width=.9\columnwidth]{Titelbild-fehlt.png}
%\caption{meine Bildunterschrift}
%label{fig:meinbild}
%\end{center}
%\end{dsafigure}



\input{LinAL_Vektoren1}
%\documentclass{article}

%\usepackage{tikz}
%\usepackage[utf8]{inputenc}
%\usepackage[ngerman]{babel}
%\usepackage{color}
%\usepackage[a4paper,lmargin={2cm},rmargin={2cm},
%tmargin={2.5cm},bmargin = {2.5cm}]{geometry}
%\usepackage{amssymb}
%\usepackage{amsmath}
%\usepackage{graphicx}
%\usepackage{multicol}
%\usepackage{amsthm}

%\newtheorem*{defi*}{Definition}
%\newtheorem*{defi2*}{Definition}

%\begin{document}
%\begin{multicols*}{2}

% \newcommand{\norm}[1]{\left\lVert#1\right\rVert}

\subsection{Norm und Skalarprodukt}

\begin{Def}
$\norm{v}$ ist eine Norm im Vektorraum $V$ falls
\begin{enumerate}
\item 
\begin{enumerate}
\item $\norm{v} \geq 0$ für alle $v \in V$
\item $\norm{v} = 0$ nur für $v = 0$
\end{enumerate}
\item $\norm{u+v} \leq \norm{u} + \norm{v}$ für alle $v \in V$ (Dreiecks-Ungleichung)
\item $\norm{\lambda v} = |\lambda| \norm{v}$ für alle $\lambda \in \mathbb{R}$ und $v \in V$ (Homogenit\"at)
\end{enumerate}
\end{Def}

In Worte gefasst ergeben sich die folgenden Eigenschaften: Eine Norm kann nie negativ sein, ist die Norm 0 bedeutet das, dass der Vektor der Nullvektor ist. Wenn man zwei Vektoren
addiert, ist die Norm des resultierenden Vektors kleiner oder gleich der Summe der Normen der beiden Ausgangsvektoren. Außerdem ist die
Norm eines Vektors, der mit einem reellen Faktor multipliziert wurde, gleich dem Produkt aus der Norm des Vektors und dem 
Betrag des Faktors.

\begin{Def}
$f:V\times V \rightarrow\mathbb{R}$ ist ein Skalarprodukt, wenn\\
\begin{enumerate}
\item $\left\langle  \alpha u + \beta v, w\right\rangle = \alpha \left\langle  u,w\right\rangle + \beta \left\langle  v,w\right\rangle $ (Bilinearität)
\item $\left\langle  u,v\right\rangle = \left\langle  v,u\right\rangle $ für alle $u,v \in V$
\item $\left\langle  v,v\right\rangle \geq 0$ mit $\left\langle  v,v\right\rangle = 0$ nur für $v = 0$
\end{enumerate}
\end{Def}
Wir können ein Skalarprodukt im euklidischen Raum $\mathbb{R}^n$ mit
\begin{equation}
\left\langle u,v\right\rangle  = u_1 v_1 + \dots + u_n v_n
\label{skal}
\end{equation}
definieren. An der ausgeschriebenen Form lässt sich überprüfen, ob es sich dabei wirklich um ein Skalarprodukt handelt. Hierbei kann  $\left\langle \alpha u + \beta v, w\right\rangle $  umgeformt werden. Durch das Ausmultiplizieren der Produkte und Sortieren der Summe in Terme mit $\alpha$ oder  $\beta$, gelangen wir zu der aus der Definition geforderten Form 
\begin{align*}
\langle \alpha u + \beta v, w\rangle &= \sum_{i=1}^n (\alpha u_i + \beta v_i)w_i \\
&= \alpha \left\langle u,w\right\rangle  + \beta \left\langle v,w\right\rangle .
\end{align*}
Die zweite Bedingung können wir mit dem Kommutativgesetz beweisen, wenn wir dieses auf die ausgeschriebene Form anwenden.
Für die dritte Bedingung bekommen wir $\sum_{i=1}^n v_i v_i = \sum_{i=1}^n v_i^2$. Da durch das Quadrieren keiner der Summanden kleiner als 0 werden kann und auch nur für $v_i = 0$ genau 0 werden kann, ist auch diese Bedingung erfüllt.

\begin{dsafigure}
\centering
\begin{tikzpicture}
\coordinate (A) at (0,0);
\coordinate (B) at (4,1.5);
\coordinate (C) at (1,3);
\draw[thick, ->] (A) -- (B) node[pos=0.5,sloped,below]{$v$};
\draw[thick, ->] (A) -- (C) node[pos=0.5,sloped, above]{$u$};
\draw[dashed, ->] (B) -- (C) node[pos=0.5,sloped, above]{$v-u$};
\draw (1,0.375) arc (21:75:1);
\node[] at (45:0.7) {$\theta$};
\end{tikzpicture}
\caption {Anwendung des Kosinussatzes}
\label{Kosinussatz}
\end{dsafigure}

Da grundsätzlich $\langle v,v\rangle  = \norm{v}^2$ gilt, kann über die Abbildung \ref{Kosinussatz} auch die Aussage getroffen werden, dass $\norm{v-u}^2 = \left\langle v-u,v-u\right\rangle $. Mit der Bilinearität des Skalarproduktes und der eben getroffenen Aussage können wir die Gleichung zu 
$\lVert v-u\rVert ^2 = \lVert v\rVert ^2 - 2\left\langle v,u\right\rangle  + \lVert u\rVert ^2$
umformen.
Zusätzlich folgt aus dem Kosinussatz, dass $\norm{v-u}^2 = \norm{u}^2 + \norm{v}^2 - 2 \norm{u} \norm{v} \cos\theta$. Setzt man beide Terme gleich, gelangt man über einige wenige Umformungen zu der Form $\left\langle u,v\right\rangle  = \norm{u} \norm{v} \cos\theta$, was ebenfalls eine Möglichkeit ist, das Skalarprodukt darzustellen.
Da sich der Betrag des Kosinus nur zwischen 0 und 1 bewegt, folgt aus dieser Gleichung direkt die Cauchy-Schwarz Ungleichung
\begin{equation}
\left\langle u,v\right\rangle  \leq \norm{u} \norm{v}.
\label{cauchy}
\end{equation}

Eine häufig genutzte Norm ist die Euklidische Norm, welche durch $\norm{v}_2 = \sqrt{v_1^2 + \dots + v_n^2}$ definiert ist und auch durch die obrige Definition des Skalarproduktes \eqref{skal} induziert wird.

Die erste und die dritte Bedingung aus der Definition für eine Norm zeigen sich aus der ausgeschriebenen Form der Euklidischen Norm. Da alle Komponenten dabei quadriert werden und die Wurzel gezogen wird, muss die erste Bedingung stimmen. Ein Faktor $\lambda$, der ebenfalls quadriert in der Wurzel steht kann aus der Summe ausgeklammert werden und als Faktor vor die Wurzel gelangen. Der Betrag ergibt sich dabei daraus, dass $\lambda$ nach dem Vorgang nur positiv sein kann, auch wenn es zuvor negativ war.

Um zu beweisen, dass die Dreiecksungleichung $\norm{u+v} \leq\norm{u} + \norm{v}$  gilt, benötigen wir die Dreiecksungleichung für reelle Zahlen und die Cauchy-Schwarz Ungleichung. Wir gehen dabei von der quadrierten Form des linken Terms aus und erhalten
\begin{align*}
\norm{u+v}^2 &= |\left\langle u+v,u+v\right\rangle|\\
&\leq |\left\langle u,u\right\rangle| + 2|\left\langle u,v\right\rangle| + |\left\langle v,v\right\rangle|\\
&\leq \norm{v}^2 + 2\norm{u} \norm{v} + \norm{v}^2\\
&= (\norm{u} + \norm{v})^2.
\end{align*}

Man sieht sofort, dass der Ausgangsterm kleiner oder gleich dem Ergebnis ist, womit bewiesen ist, dass die Euklidische Norm eine Norm ist.

%\end{multicols*}
%\end{document}

\subsection{Lineare Abbildungen}
In diesem Abschnitt werden wir die elementarste Form einer Funktion kennenlernen, die einen Vektorraum auf einen anderen Vektorraum abbildet. Diese Funktionen nennt man \textit{lineare Abbildungen}.
\begin{Def}
	\label{Def:Def_1}
Eine Funktion $f:U \stackrel{}{\rightarrow} V$ ist linear, wenn folgende Bedingungen gelten: \\ 
1. $f(u+v) = f(u)+f(v) \ \forall \ u,v \in U$\\
2. $f(\lambda v) = \lambda f(v) \ \forall \ v \in U \ \text{und} \ \lambda \in \mathbb{R}$
\end{Def}

Das besondere an solchen Abbildungen ist, dass sie von \textit{Matrizen} eindeutig definiert werden. Wie Matrizen definiert sind und wie wir sie verwenden können, lernen wir im nächsten Kapitel. Im Folgenden gehen wir davon aus, dass $U$ die Dimension $n$ hat und $V$ die Dimension $m$. 

\subsubsection{Matrix}
Eine Matrix ist eine spezielle Anordnung von $n\cdot m$ Zahlen in $n$ Zeilen und $m$ Spalten. Man kann eine Matrix aber auch als eine Zusammenfassung von $m$ Vektoren mit $n$ Komponenten ansehen, wobei ein Vektor eine Spalte bildet. Die Komponente in der $i$-ten Zeile und $j$-ten Spalte wird als $a_{ij}$ geschrieben und die ganze Spalte mit $a_j$ notiert.\\
\begin{equation*}
A= \left(
   \begin{array}{cccc}
	  a_{11} & a_{12} & ... & a_{1m}\\
		a_{21} & a_{22} & ... & a_{2m}\\
		...   & ...   & ... & ...  \\
		a_{n1} & a_{n2} & ... & a_{nm}
	 \end{array}
	 \right) = (a_1, a_2, ..., a_m)\\
\end{equation*}
\\
Da wir nun wissen was eine Matrix ist, führen wir in dem folgenden Abschnitt zwei neue Operationen ein: die \textit{Matrix-Vektor-Multiplikation} und \textit{Matrix-Matrix-Multiplikation}.

\subsubsection{Matrix-Vektor-Multiplikation}
Bei einer Matrix-Vektor-Multiplikation wird eine Matrix $A \in  \mathbb{R}^{n \times m}$ mit einem Vektor $v \in \mathbb{R}^m$ multipliziert. Dabei bildet jede Multiplikation zwischen einer Matrixzeile und einem Vektor ein Skalarprodukt. Eine Multiplikation zwischen Matrizen und Vektoren, bei denen die Anzahl nicht übereinstimmt, ist nicht definiert. Deshalb ist es wichtig, dass die Matrix so viele Spalten wie der Vektor Dimensionen hat.
\begin{align*}
	Av &= 
	\left(
   \begin{array}{cccc}
	  a_{11} & a_{12} & ... & a_{1m}\\
		a_{21} & a_{22} & ... & a_{2m}\\
		...   & ...   & ... & ...  \\
		a_{n1} & a_{n2} & ... & a_{nm}
	 \end{array}
	\right) 
	\cdot 
	\left(
	 \begin{array}{c}
	  v_1\\
		v_2\\
		...\\
		v_m
	 \end{array}
	\right)\\
	&= \left(
	 \begin{array}{c}
	 a_{11} \cdot v_1 + a_{12} \cdot v_2 + ... + a_{1m} \cdot v_m\\
	 a_{21} \cdot v_1 + a_{22} \cdot v_2 + ... + a_{2m} \cdot v_m\\
		...\\
	 a_{n1} \cdot v_1 + a_{n2} \cdot v_2 + ... + a_{nm} \cdot v_m\\
	 \end{array}
	\right)
\end{align*}

Allgemein lässt sich die Matrix-Vektor-Multiplikation folgendermaßen berechnen

\begin{equation*}
(Av)_i = \sum_{j=1}^{n}a_{ij}v_j ,
\end{equation*}

wobei $(AV)_i$ der $i$-te Eintrag des Vektors $n$ ist, der das Ergebnis der Multiplikation ist.\\
Eine Eigenschaft der linearen Abbildung ist, dass sich jede Abbildung als ein solches Matrix-Vektor-Produkt schreiben lässt. 

\subsubsection{Matrix-Matrix-Multiplikation}
Eine Matrix-Matrix-Multiplikation ist nichts anderes als mehrere Matrix-Vektor-Multiplikationen hintereinander. Haben wir zwei Matrizen $A \in \mathbb{R}^{n \times m}$ und $B \in \mathbb{R}^{m \times p}$ und berechnen $AB$, nehmen wir jede Spalte von $B$ als Vektor und multiplizieren ihn mit $A$ wie oben beschrieben. Die Vektoren, die man als Ergebnisse erhält, fasst man wieder in einer Matrix mit der Dimension $\mathbb{R}^{n \times p}$ zusammen. Allgemein können wir die Matrix-Matrix-Multiplikation folgendermaßen berechnen
\begin{equation*}
	(AB)_{il} = \sum_{n=i}^{m}a_{ij} \cdot b_{jl}
\end{equation*}

\subsubsection{Komposition}
\begin{Thm}
Eine Komposition $h(x) = (f \circ g)(x)$ aus den linearen Abbildungen $f$ und $g$, ist auch eine lineare Abbildung.
\end{Thm}

\begin{proof}
Da $f$ und $g$ linear sind gelten die Bedingungen aus Definition \ref{Def:Def_1}:
\begin{align*}
	h(u+v) &= (f \circ g)(u+v)\\
	&= f(g(u+v))\\
	&= f(g(u)+g(v))\\
	&= f(g(u)) + f(g(v))\\
	&= (f\circ g)(u) + (f\circ g)(v) = h(u) + h(v)
\end{align*}
Des Weiteren gilt:
\begin{align*}
	h(\lambda v) &= (f \circ g)(\lambda v)\\
	&=f(g(\lambda v))\\
	&= f(\lambda g(v))\\
	&=\lambda f(g(v))\\
	&=\lambda (f \circ g)(v) = \lambda h(v)
\end{align*}
Da beide Bedingungen aus Definition \ref{Def:Def_1} für $h$ erfüllt sind, ist auch eine Komposition aus zwei linearen Abbildungen linear.
\end{proof}



%\input{Kapitel-2-Kapitel�berschrift}
%\input{Kapitel-3-Kapitel�berschrift}
%usw.

%\nocite{*} % alle referenzen anzeigen, sogar wenn sie nicht im text zitiert sind
\bibliography{lit}{}
\bibliographystyle{plain}
\end{document}
